% =============================================================================
%  ERIT Benchmark Methodology — paste into Overleaf
%  Sections: Dataset Overview, Evaluation Protocol, Robustness Conditions,
%            Stratified Analysis, Metrics
% =============================================================================

% ─── 3.1 Dataset Overview ────────────────────────────────────────────────

\subsection{Dataset Overview}
\label{sec:dataset}

ERIT comprises 2{,}379 image--text pairs derived from 1{,}298 unique
ElderReact~\cite{ma2019elderreact} videos of older adults reacting to
multimedia stimuli.  Each sample consists of (i) a \emph{representative frame}
selected by a three-model ensemble (DeepFace, CLIP~ViT-B/32,
ViT-face-expression), and (ii) a \emph{spoken transcript} produced by Whisper
Large~V3 with per-token confidence scores.  Samples are annotated with one of
seven emotion labels---happiness, sadness, disgust, anger, surprise, fear, and
neutral---following the original ElderReact taxonomy.

A distinctive property of ERIT is its high degree of \emph{label ambiguity}:
74.9\% of samples carry multiple valid emotion annotations from the original
crowdsourced labels, reflecting the genuine perceptual uncertainty inherent in
older adult facial expressions (Table~\ref{tab:dataset}).

\begin{table}[t]
\centering
\caption{ERIT dataset statistics.}
\label{tab:dataset}
\begin{tabular}{lr}
\toprule
\textbf{Property} & \textbf{Value} \\
\midrule
Total samples             & 2{,}379 \\
Unique videos             & 1{,}298 \\
Emotion classes           & 7 \\
Multi-label samples       & 1{,}783 (74.9\%) \\
Single-label samples      &   596 (25.1\%) \\
Mean Whisper confidence   & 0.87 $\pm$ 0.13 \\
\midrule
\multicolumn{2}{l}{\textit{Emotion distribution (primary label)}} \\
\quad Happiness            & 730 (30.7\%) \\
\quad Surprise             & 556 (23.4\%) \\
\quad Anger                & 345 (14.5\%) \\
\quad Disgust              & 273 (11.5\%) \\
\quad Sadness              & 238 (10.0\%) \\
\quad Fear                 & 148 \phantom{0}(6.2\%) \\
\quad Neutral              &  89 \phantom{0}(3.7\%) \\
\bottomrule
\end{tabular}
\end{table}


% ─── 3.2 Evaluation Protocol ─────────────────────────────────────────────

\subsection{Evaluation Protocol}
\label{sec:protocol}

We define ERIT as a \emph{benchmark suite} rather than a single test set,
evaluating models across three orthogonal axes: modality, model regime, and
input robustness.  Detailed experimental design and results for each axis
are presented jointly in Section~\ref{sec:results}; here we summarise the
overall protocol.

\paragraph{Modality conditions.}
Every model is evaluated under three input configurations:
\begin{itemize}[nosep]
  \item \textbf{Image-only (I)}: the model receives only the representative
        frame;
  \item \textbf{Text-only (T)}: the model receives only the Whisper
        transcript (with a dummy gray image for architectures that require
        visual input);
  \item \textbf{Fusion (I+T)}: the model receives both the image and
        transcript.
\end{itemize}

\paragraph{Model regimes.}
We evaluate two families of vision--language models in two regimes:
\begin{itemize}[nosep]
  \item \textbf{Zero-shot}: LLaVA-NeXT (Mistral-7B) and Qwen2-VL (7B),
        prompted to classify emotions without task-specific training;
  \item \textbf{Fine-tuned}: LLaVA-NeXT with LoRA adapters trained on
        the ERIT training partition (1{,}169 samples).
\end{itemize}

\noindent All models are evaluated on all 2{,}379 samples under the
zero-shot regime; the fine-tuned model is evaluated on the held-out test
partition ($n\!=\!576$) only, with the dev set ($n\!=\!634$) reserved for
checkpoint selection during training.

\paragraph{Robustness conditions.}
Five perturbation conditions (grayscale, blurred faces, keyword removal,
text dropout, image dropout) are combined with a baseline condition to
yield 14 non-degenerate evaluation cells per model; see
Section~\ref{sec:robustness-results} for definitions and results.

\paragraph{Stratified analysis.}
Performance is additionally broken down by label agreement and modality
conflict; see Section~\ref{sec:stratified-results}.


% ─── 3.3 Metrics ─────────────────────────────────────────────────────────

\subsection{Metrics}
\label{sec:metrics}

For each evaluation cell we report:
\begin{itemize}[nosep]
  \item \textbf{Accuracy (Acc)}: fraction of predictions matching the
        primary ground-truth label;
  \item \textbf{Balanced Accuracy (BAcc)}: mean per-class recall,
        accounting for class imbalance;
  \item \textbf{Macro F1}: harmonic mean of per-class precision and
        recall, equally weighting all seven emotions.
\end{itemize}

\noindent
Following the multi-annotator nature of ERIT, we additionally report
\emph{multi-label accuracy}: a prediction is counted correct if it matches
\emph{any} label in the original annotation set.
Statistical significance between conditions is assessed via McNemar's test
($\alpha = 0.05$) with bootstrap 95\% confidence intervals ($B = 10{,}000$).
